\section{Механизм аудита и структура audit-логов Kubernetes}

Кластер проверяет действия, совершаемые пользователями, приложениями, использующими API Kubernetes, и самой плоскостью управления. Каждый запрос на различных этапах его выполнения создаёт событие аудита, которое затем обрабатывается в соответствии с определённой политикой и записывается в бэкэнд. Политика определяет, что именно записывается, а бэкэнды сохраняют записи. В настоящее время в качестве бэкендов используются файловая система и веб-хуки \cite{KubernetesAudit}.

Audit-логи представляют собой важный источник информации о поведении пользователей и сервисов в кластере и широко применяются для задач мониторинга безопасности, расследования инцидентов и анализа активности API. Однако из-за распределённой природы Kubernetes одно пользовательское действие может инициировать цепочку вторичных операций, выполняемых различными компонентами системы, что приводит к появлению большого числа взаимосвязанных записей в журнале. Как отмечается в работе \cite{K8NTEXT}, такие события часто оказываются разбросанными по журналу и не имеют явной связи между собой, что затрудняет реконструкцию исходного действия и анализ контекста происходящих операций.

Механизм аудита создаёт события на различных стадиях обработки запроса. Определены следующие стадии:

\begin{itemize}
    \item \texttt{RequestReceived} --- событие фиксируется сразу при получении запроса аудиторским обработчиком (до делегирования дальше по обработчикам);
    \item \texttt{ResponseStarted} --- заголовки ответа отправлены, тело ответа ещё не завершено (генерируется для долгоживущих запросов, например \texttt{watch});
    \item \texttt{ResponseComplete} --- тело ответа полностью отправлено;
    \item \texttt{Panic} --- события, генерируемые при возникновении panic в процессе обработки.
\end{itemize}

Политика содержит упорядоченный список правил: при обработке события применяется первое совпавшее правило, которое и задаёт уровень аудита для события. Доступные уровни:

\begin{itemize}
    \item \texttt{None} --- не логировать события, соответствующие правилу;
    \item \texttt{Metadata} --- логировать только метаданные (пользователь, таймштамп, ресурс, verb и т.д.), без тела запроса/ответа;
    \item \texttt{Request} --- логировать метаданные и тело запроса (не применяется к non-resource запросам);
    \item \texttt{RequestResponse} --- логировать метаданные, тело запроса и тело ответа (не применяется к non-resource запросам).
\end{itemize}